%%%%%%%%%%%%%%%%%%%%%%%%%%%%%%%%%%%%%%%%%%%%%%%%%%%%%%%%%%%%
% 5.6 MEASURE THEORY
% The Composition of Advanced Mathematics
%%%%%%%%%%%%%%%%%%%%%%%%%%%%%%%%%%%%%%%%%%%%%%%%%%%%%%%%%%%%

\section{Measure Theory}
\label{sec:measure-theory}

\prerequisite{
Real analysis, sets and functions, countability, sequences and series,
limits, continuity, basic topology of the real line, and Riemann
integration.
}

\learningobjectives{
By the end of this section, the reader should be able to:
\begin{itemize}
    \item explain why classical notions of length, area, and volume
          require generalization;
    \item work with algebras and sigma-algebras of sets;
    \item define measurable spaces and measurable functions;
    \item understand measures and their basic properties;
    \item distinguish finite, sigma-finite, and probability measures;
    \item work with Borel sets and Borel measures;
    \item understand outer measure and the construction of Lebesgue measure;
    \item apply the Carath\'eodory measurability criterion;
    \item prove basic consequences of countable additivity;
    \item use continuity from below and continuity from above;
    \item distinguish null sets from empty sets;
    \item interpret statements that hold almost everywhere;
    \item work with measurable functions and level sets;
    \item understand convergence in measure and almost-everywhere convergence;
    \item recognize the role of measure theory in integration,
          probability, functional analysis, and modern analysis.
\end{itemize}
}

%%%%%%%%%%%%%%%%%%%%%%%%%%%%%%%%%%%%%%%%%%%%%%%%%%%%%%%%%%%%
\subsection{Why Measure Theory?}
\label{subsec:why-measure-theory}
%%%%%%%%%%%%%%%%%%%%%%%%%%%%%%%%%%%%%%%%%%%%%%%%%%%%%%%%%%%%

Elementary geometry assigns lengths to intervals, areas to planar regions,
and volumes to solids.

For an interval,

\[
[a,b],
\]

its length is

\[
\boxed{
b-a.
}
\]

For a rectangle,

\[
[a,b]\times[c,d],
\]

its area is

\[
\boxed{
(b-a)(d-c).
}
\]

These ideas work well for simple geometric objects, but analysis requires
a theory capable of assigning size to far more complicated sets.

Measure theory asks:

\[
\boxed{
\text{How can the notion of size be extended consistently
to complicated sets?}
}
\]

The theory must interact correctly with countable unions, limits, and
functions.

%%%%%%%%%%%%%%%%%%%%%%%%%%%%%%%%%%%%%%%%%%%%%%%%%%%%%%%%%%%%
\subsection{The Need for Countable Operations}
\label{subsec:countable-operations-measure}
%%%%%%%%%%%%%%%%%%%%%%%%%%%%%%%%%%%%%%%%%%%%%%%%%%%%%%%%%%%%

Analysis frequently involves infinite processes.

For example,

\[
A=\bigcup_{n=1}^{\infty}A_n.
\]

A useful theory of measurable sets must therefore remain stable under
countable unions and complements.

This motivates the concept of a sigma-algebra.

%%%%%%%%%%%%%%%%%%%%%%%%%%%%%%%%%%%%%%%%%%%%%%%%%%%%%%%%%%%%
\subsection{Algebras of Sets}
\label{subsec:algebras-of-sets}
%%%%%%%%%%%%%%%%%%%%%%%%%%%%%%%%%%%%%%%%%%%%%%%%%%%%%%%%%%%%

Let \(X\) be a set.

An algebra of subsets of \(X\) is a collection \(\mathcal{A}\) satisfying:

\[
X\in\mathcal{A},
\]

\[
A\in\mathcal{A}
\Longrightarrow
A^c\in\mathcal{A},
\]

and

\[
A,B\in\mathcal{A}
\Longrightarrow
A\cup B\in\mathcal{A}.
\]

An algebra is therefore closed under finite unions, finite intersections,
and complements.

Measure theory generally requires the stronger condition of closure under
countable unions.

%%%%%%%%%%%%%%%%%%%%%%%%%%%%%%%%%%%%%%%%%%%%%%%%%%%%%%%%%%%%
\subsection{Sigma-Algebras}
\label{subsec:sigma-algebras}
%%%%%%%%%%%%%%%%%%%%%%%%%%%%%%%%%%%%%%%%%%%%%%%%%%%%%%%%%%%%

\begin{definitionbox}{Sigma-Algebra}
Let \(X\) be a set. A collection

\[
\mathcal{M}\subseteq\mathcal{P}(X)
\]

is a sigma-algebra on \(X\) if:

\[
\boxed{
X\in\mathcal{M},
}
\]

\[
\boxed{
A\in\mathcal{M}
\Longrightarrow
A^c\in\mathcal{M},
}
\]

and whenever

\[
A_1,A_2,\ldots\in\mathcal{M},
\]

then

\[
\boxed{
\bigcup_{n=1}^{\infty}A_n\in\mathcal{M}.
}
\]

The symbol \(\sigma\) is lowercase Greek \emph{sigma}. Here it indicates
closure under countable operations.
\end{definitionbox}

%%%%%%%%%%%%%%%%%%%%%%%%%%%%%%%%%%%%%%%%%%%%%%%%%%%%%%%%%%%%
\subsection{Consequences of the Sigma-Algebra Axioms}
\label{subsec:sigma-algebra-consequences}
%%%%%%%%%%%%%%%%%%%%%%%%%%%%%%%%%%%%%%%%%%%%%%%%%%%%%%%%%%%%

If \(\mathcal{M}\) is a sigma-algebra, then

\[
\boxed{
\varnothing\in\mathcal{M}.
}
\]

Indeed,

\[
\varnothing=X^c.
\]

Sigma-algebras are also closed under countable intersections because

\[
\boxed{
\bigcap_{n=1}^{\infty}A_n
=
\left(
\bigcup_{n=1}^{\infty}A_n^c
\right)^c.
}
\]

They are therefore closed under set differences as well:

\[
\boxed{
A\setminus B
=
A\cap B^c.
}
\]

%%%%%%%%%%%%%%%%%%%%%%%%%%%%%%%%%%%%%%%%%%%%%%%%%%%%%%%%%%%%
\subsection{Examples of Sigma-Algebras}
\label{subsec:examples-sigma-algebras}
%%%%%%%%%%%%%%%%%%%%%%%%%%%%%%%%%%%%%%%%%%%%%%%%%%%%%%%%%%%%

The smallest sigma-algebra on \(X\) is

\[
\boxed{
\{\varnothing,X\}.
}
\]

This is called the trivial sigma-algebra.

The largest sigma-algebra is the power set

\[
\boxed{
\mathcal{P}(X).
}
\]

If \(X\) is finite, many intermediate sigma-algebras may also exist.

%%%%%%%%%%%%%%%%%%%%%%%%%%%%%%%%%%%%%%%%%%%%%%%%%%%%%%%%%%%%
\subsection{Measurable Spaces}
\label{subsec:measurable-spaces}
%%%%%%%%%%%%%%%%%%%%%%%%%%%%%%%%%%%%%%%%%%%%%%%%%%%%%%%%%%%%

\begin{definitionbox}{Measurable Space}
A measurable space is a pair

\[
\boxed{
(X,\mathcal{M}),
}
\]

where \(X\) is a set and \(\mathcal{M}\) is a sigma-algebra on \(X\).
\end{definitionbox}

The elements of \(\mathcal{M}\) are called measurable sets.

At this stage, no numerical notion of size has yet been assigned.

%%%%%%%%%%%%%%%%%%%%%%%%%%%%%%%%%%%%%%%%%%%%%%%%%%%%%%%%%%%%
\subsection{Generated Sigma-Algebras}
\label{subsec:generated-sigma-algebras}
%%%%%%%%%%%%%%%%%%%%%%%%%%%%%%%%%%%%%%%%%%%%%%%%%%%%%%%%%%%%

Given a collection

\[
\mathcal{C}\subseteq\mathcal{P}(X),
\]

there exists a smallest sigma-algebra containing \(\mathcal{C}\).

It is denoted

\[
\boxed{
\sigma(\mathcal{C}).
}
\]

It may be defined as the intersection of all sigma-algebras containing
\(\mathcal{C}\).

Thus

\[
\boxed{
\sigma(\mathcal{C})
=
\bigcap
\{
\mathcal{M}:
\mathcal{M}\text{ is a sigma-algebra and }
\mathcal{C}\subseteq\mathcal{M}
\}.
}
\]

%%%%%%%%%%%%%%%%%%%%%%%%%%%%%%%%%%%%%%%%%%%%%%%%%%%%%%%%%%%%
\subsection{Borel Sigma-Algebra}
\label{subsec:borel-sigma-algebra}
%%%%%%%%%%%%%%%%%%%%%%%%%%%%%%%%%%%%%%%%%%%%%%%%%%%%%%%%%%%%

\begin{definitionbox}{Borel Sigma-Algebra}
The Borel sigma-algebra on \(\R\), denoted

\[
\boxed{
\mathcal{B}(\R),
}
\]

is the sigma-algebra generated by the open subsets of \(\R\).
\end{definitionbox}

Equivalently, it is generated by open intervals:

\[
\boxed{
\mathcal{B}(\R)
=
\sigma\bigl(
\{(a,b):a<b\}
\bigr).
}
\]

Its elements are called Borel sets.

Every open set and every closed set is Borel.

Countable unions and intersections of such sets are also Borel.

%%%%%%%%%%%%%%%%%%%%%%%%%%%%%%%%%%%%%%%%%%%%%%%%%%%%%%%%%%%%
\subsection{Borel Sets in Metric Spaces}
\label{subsec:borel-metric-spaces}
%%%%%%%%%%%%%%%%%%%%%%%%%%%%%%%%%%%%%%%%%%%%%%%%%%%%%%%%%%%%

More generally, if \(X\) is a topological space, then

\[
\boxed{
\mathcal{B}(X)
}
\]

denotes the sigma-algebra generated by the open subsets of \(X\).

This creates an important bridge between topology and measure theory.

Topology determines which sets are open.

The Borel construction then closes those sets under countable
set-theoretic operations.

%%%%%%%%%%%%%%%%%%%%%%%%%%%%%%%%%%%%%%%%%%%%%%%%%%%%%%%%%%%%
\subsection{Measures}
\label{subsec:measures}
%%%%%%%%%%%%%%%%%%%%%%%%%%%%%%%%%%%%%%%%%%%%%%%%%%%%%%%%%%%%

\begin{definitionbox}{Measure}
Let

\[
(X,\mathcal{M})
\]

be a measurable space.

A measure is a function

\[
\boxed{
\mu:\mathcal{M}\to[0,\infty]
}
\]

satisfying

\[
\boxed{
\mu(\varnothing)=0
}
\]

and countable additivity:

if \(A_1,A_2,\ldots\in\mathcal{M}\) are pairwise disjoint, then

\[
\boxed{
\mu\left(
\bigcup_{n=1}^{\infty}A_n
\right)
=
\sum_{n=1}^{\infty}\mu(A_n).
}
\]

The symbol \(\mu\) is lowercase Greek \emph{mu}; it commonly denotes a
measure.
\end{definitionbox}

The codomain includes

\[
+\infty,
\]

so measurable sets are permitted to have infinite measure.

%%%%%%%%%%%%%%%%%%%%%%%%%%%%%%%%%%%%%%%%%%%%%%%%%%%%%%%%%%%%
\subsection{Measure Spaces}
\label{subsec:measure-spaces}
%%%%%%%%%%%%%%%%%%%%%%%%%%%%%%%%%%%%%%%%%%%%%%%%%%%%%%%%%%%%

\begin{definitionbox}{Measure Space}
A measure space is a triple

\[
\boxed{
(X,\mathcal{M},\mu),
}
\]

where

\[
(X,\mathcal{M})
\]

is a measurable space and \(\mu\) is a measure on \(\mathcal{M}\).
\end{definitionbox}

This structure forms the basic setting for measure-theoretic integration.

%%%%%%%%%%%%%%%%%%%%%%%%%%%%%%%%%%%%%%%%%%%%%%%%%%%%%%%%%%%%
\subsection{Finite Measures}
\label{subsec:finite-measures}
%%%%%%%%%%%%%%%%%%%%%%%%%%%%%%%%%%%%%%%%%%%%%%%%%%%%%%%%%%%%

A measure \(\mu\) on \(X\) is finite if

\[
\boxed{
\mu(X)<\infty.
}
\]

If the whole space has finite measure, then every measurable subset also
has finite measure.

%%%%%%%%%%%%%%%%%%%%%%%%%%%%%%%%%%%%%%%%%%%%%%%%%%%%%%%%%%%%
\subsection{Sigma-Finite Measures}
\label{subsec:sigma-finite-measures}
%%%%%%%%%%%%%%%%%%%%%%%%%%%%%%%%%%%%%%%%%%%%%%%%%%%%%%%%%%%%

\begin{definitionbox}{Sigma-Finite Measure}
A measure space is sigma-finite if there exist measurable sets

\[
E_1,E_2,\ldots
\]

such that

\[
\boxed{
X=
\bigcup_{n=1}^{\infty}E_n
}
\]

and

\[
\boxed{
\mu(E_n)<\infty
}
\]

for every \(n\).
\end{definitionbox}

Lebesgue measure on \(\R\) is not finite because

\[
m(\R)=\infty,
\]

but it is sigma-finite since

\[
\R
=
\bigcup_{n=1}^{\infty}[-n,n]
\]

and

\[
m([-n,n])=2n<\infty.
\]

%%%%%%%%%%%%%%%%%%%%%%%%%%%%%%%%%%%%%%%%%%%%%%%%%%%%%%%%%%%%
\subsection{Probability Measures}
\label{subsec:probability-measures}
%%%%%%%%%%%%%%%%%%%%%%%%%%%%%%%%%%%%%%%%%%%%%%%%%%%%%%%%%%%%

A probability measure is a measure satisfying

\[
\boxed{
\mu(X)=1.
}
\]

In probability theory, the notation

\[
\mathbb{P}
\]

is commonly used instead of \(\mu\).

Thus a probability space is written

\[
\boxed{
(\Omega,\mathcal{F},\mathbb{P}).
}
\]

The uppercase Greek letter

\[
\Omega
\]

is \emph{Omega} and usually denotes the sample space.

%%%%%%%%%%%%%%%%%%%%%%%%%%%%%%%%%%%%%%%%%%%%%%%%%%%%%%%%%%%%
\subsection{Counting Measure}
\label{subsec:counting-measure}
%%%%%%%%%%%%%%%%%%%%%%%%%%%%%%%%%%%%%%%%%%%%%%%%%%%%%%%%%%%%

For any set \(X\), define

\[
\boxed{
\mu(A)
=
\#A
}
\]

when \(A\) is finite and

\[
\mu(A)=\infty
\]

when \(A\) is infinite.

This is called counting measure.

For example,

\[
\mu(\{a,b,c\})=3.
\]

%%%%%%%%%%%%%%%%%%%%%%%%%%%%%%%%%%%%%%%%%%%%%%%%%%%%%%%%%%%%
\subsection{Dirac Measure}
\label{subsec:dirac-measure}
%%%%%%%%%%%%%%%%%%%%%%%%%%%%%%%%%%%%%%%%%%%%%%%%%%%%%%%%%%%%

Fix a point

\[
x_0\in X.
\]

The Dirac measure at \(x_0\) is defined by

\[
\boxed{
\delta_{x_0}(A)
=
\begin{cases}
1, & x_0\in A,\\
0, & x_0\notin A.
\end{cases}
}
\]

Here \(\delta\) is lowercase Greek \emph{delta}.

The Dirac measure concentrates all of its mass at one point.

%%%%%%%%%%%%%%%%%%%%%%%%%%%%%%%%%%%%%%%%%%%%%%%%%%%%%%%%%%%%
\subsection{Basic Properties of Measures}
\label{subsec:basic-properties-measures}
%%%%%%%%%%%%%%%%%%%%%%%%%%%%%%%%%%%%%%%%%%%%%%%%%%%%%%%%%%%%

Countable additivity implies several important properties.

If

\[
A\subseteq B,
\]

then

\[
\boxed{
\mu(A)\leq\mu(B).
}
\]

This is monotonicity.

If \(A\subseteq B\) and \(\mu(A)<\infty\), then

\[
\boxed{
\mu(B\setminus A)
=
\mu(B)-\mu(A).
}
\]

For arbitrary measurable \(A\) and \(B\),

\[
\boxed{
\mu(A\cup B)
+
\mu(A\cap B)
=
\mu(A)+\mu(B)
}
\]

whenever the quantities are interpreted without an undefined
\(\infty-\infty\) expression.

%%%%%%%%%%%%%%%%%%%%%%%%%%%%%%%%%%%%%%%%%%%%%%%%%%%%%%%%%%%%
\subsection{Countable Subadditivity}
\label{subsec:countable-subadditivity}
%%%%%%%%%%%%%%%%%%%%%%%%%%%%%%%%%%%%%%%%%%%%%%%%%%%%%%%%%%%%

For measurable sets

\[
A_1,A_2,\ldots,
\]

not necessarily disjoint,

\[
\boxed{
\mu\left(
\bigcup_{n=1}^{\infty}A_n
\right)
\leq
\sum_{n=1}^{\infty}\mu(A_n).
}
\]

This property is called countable subadditivity.

Equality is guaranteed when the sets are pairwise disjoint.

%%%%%%%%%%%%%%%%%%%%%%%%%%%%%%%%%%%%%%%%%%%%%%%%%%%%%%%%%%%%
\subsection{Continuity from Below}
\label{subsec:continuity-from-below-measure}
%%%%%%%%%%%%%%%%%%%%%%%%%%%%%%%%%%%%%%%%%%%%%%%%%%%%%%%%%%%%

\begin{theorem}[Continuity from Below]
Suppose

\[
A_1\subseteq A_2\subseteq A_3\subseteq\cdots.
\]

Then

\[
\boxed{
\mu\left(
\bigcup_{n=1}^{\infty}A_n
\right)
=
\lim_{n\to\infty}\mu(A_n).
}
\]
\end{theorem}

This is often written

\[
A_n\uparrow A
\]

to indicate that the sets increase to \(A\).

Then

\[
\boxed{
\mu(A_n)\to\mu(A).
}
\]

%%%%%%%%%%%%%%%%%%%%%%%%%%%%%%%%%%%%%%%%%%%%%%%%%%%%%%%%%%%%
\subsection{Continuity from Above}
\label{subsec:continuity-from-above-measure}
%%%%%%%%%%%%%%%%%%%%%%%%%%%%%%%%%%%%%%%%%%%%%%%%%%%%%%%%%%%%

\begin{theorem}[Continuity from Above]
Suppose

\[
A_1\supseteq A_2\supseteq A_3\supseteq\cdots
\]

and

\[
\mu(A_1)<\infty.
\]

Then

\[
\boxed{
\mu\left(
\bigcap_{n=1}^{\infty}A_n
\right)
=
\lim_{n\to\infty}\mu(A_n).
}
\]
\end{theorem}

The finiteness assumption is important.

The notation

\[
A_n\downarrow A
\]

means that the sets decrease to \(A\).

%%%%%%%%%%%%%%%%%%%%%%%%%%%%%%%%%%%%%%%%%%%%%%%%%%%%%%%%%%%%
\subsection{Worked Example: Continuity from Below}
\label{subsec:worked-continuity-below}
%%%%%%%%%%%%%%%%%%%%%%%%%%%%%%%%%%%%%%%%%%%%%%%%%%%%%%%%%%%%

\begin{workedexamplebox}{Increasing Intervals}

Let

\[
A_n=
\left[0,1-\frac1n\right]
\]

for \(n\geq1\).

Then

\[
A_1\subseteq A_2\subseteq\cdots
\]

and

\[
\bigcup_{n=1}^{\infty}A_n
=
[0,1).
\]

Using ordinary length,

\[
m(A_n)
=
1-\frac1n.
\]

Therefore

\[
\lim_{n\to\infty}m(A_n)
=
1.
\]

By continuity from below,

\[
\begin{answerbox}
\[
\boxed{
m([0,1))=1.
}
\]
\end{answerbox}
\end{workedexamplebox}

%%%%%%%%%%%%%%%%%%%%%%%%%%%%%%%%%%%%%%%%%%%%%%%%%%%%%%%%%%%%
\subsection{Outer Measure}
\label{subsec:outer-measure}
%%%%%%%%%%%%%%%%%%%%%%%%%%%%%%%%%%%%%%%%%%%%%%%%%%%%%%%%%%%%

A central problem is how to construct a measure rather than merely assume
one exists.

For subsets of \(\R\), begin by covering an arbitrary set \(E\) with
countably many intervals.

The Lebesgue outer measure is defined by

\[
\boxed{
m^*(E)
=
\inf
\left\{
\sum_{n=1}^{\infty}|I_n|:
E\subseteq
\bigcup_{n=1}^{\infty}I_n
\right\}.
}
\]

Here

\[
|I_n|
\]

denotes the length of the interval \(I_n\).

The symbol

\[
m^*
\]

is read \emph{\(m\)-star}.

%%%%%%%%%%%%%%%%%%%%%%%%%%%%%%%%%%%%%%%%%%%%%%%%%%%%%%%%%%%%
\subsection{Properties of Outer Measure}
\label{subsec:properties-outer-measure}
%%%%%%%%%%%%%%%%%%%%%%%%%%%%%%%%%%%%%%%%%%%%%%%%%%%%%%%%%%%%

An outer measure \(\mu^*\) satisfies

\[
\boxed{
\mu^*(\varnothing)=0,
}
\]

monotonicity,

\[
A\subseteq B
\Longrightarrow
\boxed{
\mu^*(A)\leq\mu^*(B),
}
\]

and countable subadditivity,

\[
\boxed{
\mu^*\left(
\bigcup_{n=1}^{\infty}A_n
\right)
\leq
\sum_{n=1}^{\infty}\mu^*(A_n).
}
\]

Unlike a measure, an outer measure is typically defined on every subset of
the underlying space.

%%%%%%%%%%%%%%%%%%%%%%%%%%%%%%%%%%%%%%%%%%%%%%%%%%%%%%%%%%%%
\subsection{Carath\'eodory Measurability}
\label{subsec:caratheodory-measurability}
%%%%%%%%%%%%%%%%%%%%%%%%%%%%%%%%%%%%%%%%%%%%%%%%%%%%%%%%%%%%

\begin{definitionbox}{Carath\'eodory Measurable Set}
Let \(\mu^*\) be an outer measure on \(X\).

A set \(E\subseteq X\) is Carath\'eodory measurable if for every
\(A\subseteq X\),

\[
\boxed{
\mu^*(A)
=
\mu^*(A\cap E)
+
\mu^*(A\setminus E).
}
\]
\end{definitionbox}

The criterion says that \(E\) divides every set \(A\) cleanly with respect
to outer measure.

%%%%%%%%%%%%%%%%%%%%%%%%%%%%%%%%%%%%%%%%%%%%%%%%%%%%%%%%%%%%
\subsection{Carath\'eodory's Construction}
\label{subsec:caratheodory-construction}
%%%%%%%%%%%%%%%%%%%%%%%%%%%%%%%%%%%%%%%%%%%%%%%%%%%%%%%%%%%%

\begin{theorem}[Carath\'eodory's Theorem]
The collection of Carath\'eodory measurable sets for an outer measure
\(\mu^*\) forms a sigma-algebra.

The restriction of \(\mu^*\) to this sigma-algebra is a measure.
\end{theorem}

This provides a general method:

\[
\boxed{
\text{outer measure}
\longrightarrow
\text{measurable sets}
\longrightarrow
\text{measure}.
}
\]

Lebesgue measure is constructed using this principle.

%%%%%%%%%%%%%%%%%%%%%%%%%%%%%%%%%%%%%%%%%%%%%%%%%%%%%%%%%%%%
\subsection{Lebesgue Measure}
\label{subsec:lebesgue-measure}
%%%%%%%%%%%%%%%%%%%%%%%%%%%%%%%%%%%%%%%%%%%%%%%%%%%%%%%%%%%%

Lebesgue measure on \(\R\) is conventionally denoted

\[
\boxed{
m.
}
\]

It extends ordinary interval length:

\[
\boxed{
m([a,b])=b-a.
}
\]

The same value holds for

\[
(a,b),
\qquad
[a,b),
\qquad
(a,b].
\]

Endpoints have measure zero.

%%%%%%%%%%%%%%%%%%%%%%%%%%%%%%%%%%%%%%%%%%%%%%%%%%%%%%%%%%%%
\subsection{Translation Invariance}
\label{subsec:translation-invariance-measure}
%%%%%%%%%%%%%%%%%%%%%%%%%%%%%%%%%%%%%%%%%%%%%%%%%%%%%%%%%%%%

Lebesgue measure is translation invariant.

For a measurable set \(E\subseteq\R\) and \(t\in\R\),

\[
E+t
=
\{x+t:x\in E\}.
\]

Then

\[
\boxed{
m(E+t)=m(E).
}
\]

Thus moving a set does not change its Lebesgue measure.

%%%%%%%%%%%%%%%%%%%%%%%%%%%%%%%%%%%%%%%%%%%%%%%%%%%%%%%%%%%%
\subsection{Scaling of Lebesgue Measure}
\label{subsec:scaling-lebesgue-measure}
%%%%%%%%%%%%%%%%%%%%%%%%%%%%%%%%%%%%%%%%%%%%%%%%%%%%%%%%%%%%

For measurable

\[
E\subseteq\R
\]

and \(c\in\R\),

\[
cE
=
\{cx:x\in E\}.
\]

Then

\[
\boxed{
m(cE)
=
|c|\,m(E).
}
\]

In \(\R^n\), scalar dilation gives

\[
\boxed{
m_n(cE)
=
|c|^n m_n(E).
}
\]

This generalizes familiar geometric scaling laws.

%%%%%%%%%%%%%%%%%%%%%%%%%%%%%%%%%%%%%%%%%%%%%%%%%%%%%%%%%%%%
\subsection{Null Sets}
\label{subsec:null-sets}
%%%%%%%%%%%%%%%%%%%%%%%%%%%%%%%%%%%%%%%%%%%%%%%%%%%%%%%%%%%%

\begin{definitionbox}{Null Set}
A measurable set \(N\) is a null set if

\[
\boxed{
\mu(N)=0.
}
\]
\end{definitionbox}

A null set need not be empty.

For Lebesgue measure,

\[
\boxed{
m(\{x\})=0
}
\]

for every \(x\in\R\).

Therefore every finite subset of \(\R\) has measure zero.

%%%%%%%%%%%%%%%%%%%%%%%%%%%%%%%%%%%%%%%%%%%%%%%%%%%%%%%%%%%%
\subsection{Countable Sets Have Measure Zero}
\label{subsec:countable-null-sets}
%%%%%%%%%%%%%%%%%%%%%%%%%%%%%%%%%%%%%%%%%%%%%%%%%%%%%%%%%%%%

Let

\[
A=\{a_1,a_2,\ldots\}
\]

be countable.

Then

\[
A=
\bigcup_{n=1}^{\infty}\{a_n\}.
\]

By countable subadditivity,

\[
m(A)
\leq
\sum_{n=1}^{\infty}m(\{a_n\})
=
0.
\]

Therefore

\[
\boxed{
m(A)=0.
}
\]

In particular,

\[
\boxed{
m(\Q)=0.
}
\]

The rational numbers are dense in \(\R\) yet have Lebesgue measure zero.

%%%%%%%%%%%%%%%%%%%%%%%%%%%%%%%%%%%%%%%%%%%%%%%%%%%%%%%%%%%%
\subsection{Uncountable Null Sets}
\label{subsec:uncountable-null-sets}
%%%%%%%%%%%%%%%%%%%%%%%%%%%%%%%%%%%%%%%%%%%%%%%%%%%%%%%%%%%%

Measure zero does not imply countability.

The standard middle-thirds Cantor set is uncountable but has Lebesgue
measure zero.

Thus cardinality and measure describe fundamentally different notions of
size.

A set may be:

\[
\boxed{
\text{uncountably infinite}
}
\]

while simultaneously having

\[
\boxed{
\text{measure zero}.
}
\]

%%%%%%%%%%%%%%%%%%%%%%%%%%%%%%%%%%%%%%%%%%%%%%%%%%%%%%%%%%%%
\subsection{Almost Everywhere}
\label{subsec:almost-everywhere}
%%%%%%%%%%%%%%%%%%%%%%%%%%%%%%%%%%%%%%%%%%%%%%%%%%%%%%%%%%%%

A property is said to hold almost everywhere if the set where it fails has
measure zero.

The standard abbreviation is

\[
\boxed{
\text{a.e.}
}
\]

For example,

\[
f=g
\quad
\text{a.e.}
\]

means

\[
\boxed{
\mu(
\{x:f(x)\neq g(x)\}
)
=
0.
}
\]

Measure theory often treats functions that differ only on a null set as
essentially equivalent.

%%%%%%%%%%%%%%%%%%%%%%%%%%%%%%%%%%%%%%%%%%%%%%%%%%%%%%%%%%%%
\subsection{Complete Measures}
\label{subsec:complete-measures}
%%%%%%%%%%%%%%%%%%%%%%%%%%%%%%%%%%%%%%%%%%%%%%%%%%%%%%%%%%%%

\begin{definitionbox}{Complete Measure}
A measure space

\[
(X,\mathcal{M},\mu)
\]

is complete if every subset of every measurable null set is measurable.
\end{definitionbox}

Lebesgue measure is complete.

The Borel measure obtained by restricting Lebesgue measure to Borel sets is
not complete.

Completing it adds all subsets of Borel null sets.

%%%%%%%%%%%%%%%%%%%%%%%%%%%%%%%%%%%%%%%%%%%%%%%%%%%%%%%%%%%%
\subsection{Lebesgue Measurable Versus Borel}
\label{subsec:lebesgue-vs-borel}
%%%%%%%%%%%%%%%%%%%%%%%%%%%%%%%%%%%%%%%%%%%%%%%%%%%%%%%%%%%%

Every Borel subset of \(\R\) is Lebesgue measurable:

\[
\boxed{
\mathcal{B}(\R)
\subseteq
\mathcal{L}(\R),
}
\]

where

\[
\mathcal{L}(\R)
\]

denotes the Lebesgue sigma-algebra.

The inclusion is strict:

\[
\boxed{
\mathcal{B}(\R)
\subsetneq
\mathcal{L}(\R).
}
\]

Thus some Lebesgue measurable sets are not Borel sets.

%%%%%%%%%%%%%%%%%%%%%%%%%%%%%%%%%%%%%%%%%%%%%%%%%%%%%%%%%%%%
\subsection{Nonmeasurable Sets}
\label{subsec:nonmeasurable-sets}
%%%%%%%%%%%%%%%%%%%%%%%%%%%%%%%%%%%%%%%%%%%%%%%%%%%%%%%%%%%%

Not every subset of \(\R\) can be assigned a translation-invariant,
countably additive measure extending ordinary length.

A classical example is a Vitali set.

Its construction uses the Axiom of Choice.

The existence of nonmeasurable sets explains why Lebesgue measure cannot
simply be defined on the entire power set

\[
\mathcal{P}(\R)
\]

while retaining all desired properties.

%%%%%%%%%%%%%%%%%%%%%%%%%%%%%%%%%%%%%%%%%%%%%%%%%%%%%%%%%%%%
\subsection{Measurable Functions}
\label{subsec:measurable-functions}
%%%%%%%%%%%%%%%%%%%%%%%%%%%%%%%%%%%%%%%%%%%%%%%%%%%%%%%%%%%%

\begin{definitionbox}{Measurable Function}
Let

\[
(X,\mathcal{M})
\]

and

\[
(Y,\mathcal{N})
\]

be measurable spaces.

A function

\[
f:X\to Y
\]

is measurable if

\[
\boxed{
B\in\mathcal{N}
\Longrightarrow
f^{-1}(B)\in\mathcal{M}.
}
\]
\end{definitionbox}

This resembles the topological definition of continuity.

For continuity, inverse images of open sets must be open.

For measurability, inverse images of measurable sets must be measurable.

%%%%%%%%%%%%%%%%%%%%%%%%%%%%%%%%%%%%%%%%%%%%%%%%%%%%%%%%%%%%
\subsection{Real-Valued Measurable Functions}
\label{subsec:real-valued-measurable-functions}
%%%%%%%%%%%%%%%%%%%%%%%%%%%%%%%%%%%%%%%%%%%%%%%%%%%%%%%%%%%%

For a real-valued function

\[
f:X\to\R,
\]

it is enough to check certain families of sets.

For example, \(f\) is measurable if

\[
\boxed{
\{x\in X:f(x)>a\}\in\mathcal{M}
}
\]

for every \(a\in\R\).

Equivalent tests use sets of the forms

\[
\{f\geq a\},
\qquad
\{f<a\},
\qquad
\{f\leq a\}.
\]

%%%%%%%%%%%%%%%%%%%%%%%%%%%%%%%%%%%%%%%%%%%%%%%%%%%%%%%%%%%%
\subsection{Continuous Functions Are Borel Measurable}
\label{subsec:continuous-borel-measurable}
%%%%%%%%%%%%%%%%%%%%%%%%%%%%%%%%%%%%%%%%%%%%%%%%%%%%%%%%%%%%

If

\[
f:X\to Y
\]

is continuous between topological spaces, then \(f\) is measurable with
respect to their Borel sigma-algebras.

Indeed, inverse images of open sets are open.

Since the Borel sigma-algebra is generated by open sets, continuity implies
Borel measurability.

Thus

\[
\boxed{
\text{continuous}
\Longrightarrow
\text{Borel measurable}.
}
\]

The converse is false.

%%%%%%%%%%%%%%%%%%%%%%%%%%%%%%%%%%%%%%%%%%%%%%%%%%%%%%%%%%%%
\subsection{Indicator Functions}
\label{subsec:indicator-functions}
%%%%%%%%%%%%%%%%%%%%%%%%%%%%%%%%%%%%%%%%%%%%%%%%%%%%%%%%%%%%

For a set \(A\subseteq X\), its indicator function is

\[
\boxed{
\mathbf{1}_A(x)
=
\begin{cases}
1, & x\in A,\\
0, & x\notin A.
\end{cases}
}
\]

It is also commonly written

\[
\chi_A.
\]

The indicator function is measurable if and only if \(A\) is measurable.

Indicator functions are fundamental building blocks for integration.

%%%%%%%%%%%%%%%%%%%%%%%%%%%%%%%%%%%%%%%%%%%%%%%%%%%%%%%%%%%%
\subsection{Simple Functions}
\label{subsec:simple-functions}
%%%%%%%%%%%%%%%%%%%%%%%%%%%%%%%%%%%%%%%%%%%%%%%%%%%%%%%%%%%%

A measurable simple function takes only finitely many values.

It can be written

\[
\boxed{
\phi
=
\sum_{k=1}^{n}
a_k\mathbf{1}_{E_k},
}
\]

where each \(E_k\) is measurable.

The lowercase Greek letter

\[
\phi
\]

is \emph{phi}.

Simple functions form the first stage in constructing the Lebesgue
integral.

%%%%%%%%%%%%%%%%%%%%%%%%%%%%%%%%%%%%%%%%%%%%%%%%%%%%%%%%%%%%
\subsection{Algebra of Measurable Functions}
\label{subsec:algebra-measurable-functions}
%%%%%%%%%%%%%%%%%%%%%%%%%%%%%%%%%%%%%%%%%%%%%%%%%%%%%%%%%%%%

If \(f\) and \(g\) are real-valued measurable functions, then under their
natural domains,

\[
\boxed{
f+g,
\quad
f-g,
\quad
fg
}
\]

are measurable.

When the quotient is defined,

\[
\boxed{
\frac{f}{g}
}
\]

is measurable.

Also,

\[
\boxed{
|f|
}
\]

is measurable.

%%%%%%%%%%%%%%%%%%%%%%%%%%%%%%%%%%%%%%%%%%%%%%%%%%%%%%%%%%%%
\subsection{Limits of Measurable Functions}
\label{subsec:limits-measurable-functions}
%%%%%%%%%%%%%%%%%%%%%%%%%%%%%%%%%%%%%%%%%%%%%%%%%%%%%%%%%%%%

If

\[
f_n:X\to\overline{\R}
\]

are measurable, then the functions

\[
\boxed{
\sup_n f_n,
\qquad
\inf_n f_n,
}
\]

\[
\boxed{
\limsup_{n\to\infty}f_n,
\qquad
\liminf_{n\to\infty}f_n
}
\]

are measurable whenever interpreted as extended-real-valued functions.

If

\[
f_n(x)\to f(x)
\]

pointwise, then \(f\) is measurable.

Thus measurability behaves well under pointwise limits.

%%%%%%%%%%%%%%%%%%%%%%%%%%%%%%%%%%%%%%%%%%%%%%%%%%%%%%%%%%%%
\subsection{Almost-Everywhere Convergence}
\label{subsec:almost-everywhere-convergence}
%%%%%%%%%%%%%%%%%%%%%%%%%%%%%%%%%%%%%%%%%%%%%%%%%%%%%%%%%%%%

A sequence \((f_n)\) converges to \(f\) almost everywhere if

\[
\boxed{
f_n(x)\to f(x)
}
\]

for all \(x\) outside a null set.

Equivalently,

\[
\boxed{
\mu\left(
\left\{
x:
f_n(x)\not\to f(x)
\right\}
\right)
=
0.
}
\]

This is weaker than pointwise convergence everywhere because exceptional
null sets are permitted.

%%%%%%%%%%%%%%%%%%%%%%%%%%%%%%%%%%%%%%%%%%%%%%%%%%%%%%%%%%%%
\subsection{Convergence in Measure}
\label{subsec:convergence-in-measure}
%%%%%%%%%%%%%%%%%%%%%%%%%%%%%%%%%%%%%%%%%%%%%%%%%%%%%%%%%%%%

\begin{definitionbox}{Convergence in Measure}
A sequence \((f_n)\) converges in measure to \(f\) if for every
\(\varepsilon>0\),

\[
\boxed{
\mu\left(
\left\{
x:
|f_n(x)-f(x)|>\varepsilon
\right\}
\right)
\to0.
}
\]
\end{definitionbox}

Instead of requiring pointwise closeness, convergence in measure requires
the set on which the functions differ significantly to become small in
measure.

%%%%%%%%%%%%%%%%%%%%%%%%%%%%%%%%%%%%%%%%%%%%%%%%%%%%%%%%%%%%
\subsection{Relations Between Modes of Convergence}
\label{subsec:relations-measure-convergence}
%%%%%%%%%%%%%%%%%%%%%%%%%%%%%%%%%%%%%%%%%%%%%%%%%%%%%%%%%%%%

On a finite measure space,

\[
\boxed{
f_n\to f
\text{ a.e.}
\Longrightarrow
f_n\to f
\text{ in measure}.
}
\]

More generally, on arbitrary measure spaces this implication can fail
without additional hypotheses.

Conversely,

\[
\boxed{
\text{convergence in measure}
\not\Longrightarrow
\text{a.e. convergence}
}
\]

in general.

However, convergence in measure implies that some subsequence converges
almost everywhere under standard hypotheses.

These distinctions become important in Lebesgue integration.

%%%%%%%%%%%%%%%%%%%%%%%%%%%%%%%%%%%%%%%%%%%%%%%%%%%%%%%%%%%%
\subsection{Worked Example: Convergence in Measure}
\label{subsec:worked-convergence-measure}
%%%%%%%%%%%%%%%%%%%%%%%%%%%%%%%%%%%%%%%%%%%%%%%%%%%%%%%%%%%%

\begin{workedexamplebox}{Shrinking Indicator Functions}

On

\[
[0,1]
\]

with Lebesgue measure, define

\[
f_n
=
\mathbf{1}_{[0,1/n]}.
\]

For every fixed

\[
x>0,
\]

we eventually have

\[
x>\frac1n,
\]

so

\[
f_n(x)=0.
\]

At \(x=0\),

\[
f_n(0)=1
\]

for every \(n\).

Thus

\[
f_n(x)\to0
\]

almost everywhere.

For

\[
0<\varepsilon<1,
\]

the set where

\[
|f_n(x)|>\varepsilon
\]

is

\[
[0,1/n].
\]

Therefore

\[
m\left(
\{|f_n|>\varepsilon\}
\right)
=
\frac1n
\to0.
\]

\begin{answerbox}
\[
\boxed{
f_n\to0
\text{ almost everywhere and in measure}.
}
\]
\end{answerbox}
\end{workedexamplebox}

%%%%%%%%%%%%%%%%%%%%%%%%%%%%%%%%%%%%%%%%%%%%%%%%%%%%%%%%%%%%
\subsection{Measure-Theoretic Limits of Sets}
\label{subsec:limsup-liminf-sets}
%%%%%%%%%%%%%%%%%%%%%%%%%%%%%%%%%%%%%%%%%%%%%%%%%%%%%%%%%%%%

For a sequence of sets \((A_n)\), define

\[
\boxed{
\limsup_{n\to\infty}A_n
=
\bigcap_{n=1}^{\infty}
\bigcup_{k=n}^{\infty}A_k
}
\]

and

\[
\boxed{
\liminf_{n\to\infty}A_n
=
\bigcup_{n=1}^{\infty}
\bigcap_{k=n}^{\infty}A_k.
}
\]

A point belongs to

\[
\limsup A_n
\]

if it belongs to infinitely many \(A_n\).

A point belongs to

\[
\liminf A_n
\]

if it belongs to all sufficiently large \(A_n\).

%%%%%%%%%%%%%%%%%%%%%%%%%%%%%%%%%%%%%%%%%%%%%%%%%%%%%%%%%%%%
\subsection{Borel--Cantelli Preview}
\label{subsec:borel-cantelli-preview}
%%%%%%%%%%%%%%%%%%%%%%%%%%%%%%%%%%%%%%%%%%%%%%%%%%%%%%%%%%%%

A fundamental result connecting measure and repeated events is the
Borel--Cantelli lemma.

\begin{theorem}[First Borel--Cantelli Lemma]
If

\[
\sum_{n=1}^{\infty}\mu(A_n)<\infty,
\]

then

\[
\boxed{
\mu\left(
\limsup_{n\to\infty}A_n
\right)=0.
}
\]
\end{theorem}

Thus if the total sum of the measures is finite, almost no point belongs to
infinitely many of the sets.

This result becomes especially important in probability theory.

%%%%%%%%%%%%%%%%%%%%%%%%%%%%%%%%%%%%%%%%%%%%%%%%%%%%%%%%%%%%
\subsection{Product Measures Preview}
\label{subsec:product-measures-preview}
%%%%%%%%%%%%%%%%%%%%%%%%%%%%%%%%%%%%%%%%%%%%%%%%%%%%%%%%%%%%

Given measure spaces

\[
(X,\mathcal{M},\mu)
\]

and

\[
(Y,\mathcal{N},\nu),
\]

one seeks a measure on

\[
X\times Y
\]

satisfying

\[
\boxed{
(\mu\times\nu)(A\times B)
=
\mu(A)\nu(B).
}
\]

The measure

\[
\mu\times\nu
\]

is called a product measure.

Product measures provide the foundation for multidimensional integration
and the Fubini--Tonelli theorems.

%%%%%%%%%%%%%%%%%%%%%%%%%%%%%%%%%%%%%%%%%%%%%%%%%%%%%%%%%%%%
\subsection{Lebesgue Measure in Higher Dimensions}
\label{subsec:lebesgue-higher-dimensions}
%%%%%%%%%%%%%%%%%%%%%%%%%%%%%%%%%%%%%%%%%%%%%%%%%%%%%%%%%%%%

Lebesgue measure extends naturally to

\[
\R^n.
\]

It is commonly denoted

\[
m_n.
\]

For a rectangular box

\[
R=
\prod_{k=1}^{n}[a_k,b_k],
\]

its measure is

\[
\boxed{
m_n(R)
=
\prod_{k=1}^{n}(b_k-a_k).
}
\]

Thus Lebesgue measure generalizes length, area, volume, and
higher-dimensional volume.

%%%%%%%%%%%%%%%%%%%%%%%%%%%%%%%%%%%%%%%%%%%%%%%%%%%%%%%%%%%%
\subsection{Regularity of Lebesgue Measure}
\label{subsec:regularity-lebesgue-measure}
%%%%%%%%%%%%%%%%%%%%%%%%%%%%%%%%%%%%%%%%%%%%%%%%%%%%%%%%%%%%

Lebesgue measure can approximate measurable sets using geometrically
simpler sets.

For a Lebesgue measurable set \(E\),

\[
\boxed{
m(E)
=
\inf
\{
m(U):
E\subseteq U,\;
U\text{ open}
\}.
}
\]

For measurable sets of finite measure,

\[
\boxed{
m(E)
=
\sup
\{
m(K):
K\subseteq E,\;
K\text{ compact}
\}.
}
\]

These are forms of outer and inner regularity.

%%%%%%%%%%%%%%%%%%%%%%%%%%%%%%%%%%%%%%%%%%%%%%%%%%%%%%%%%%%%
\subsection{Measure and Topology}
\label{subsec:measure-and-topology}
%%%%%%%%%%%%%%%%%%%%%%%%%%%%%%%%%%%%%%%%%%%%%%%%%%%%%%%%%%%%

Topology and measure theory describe different kinds of structure.

Topology studies concepts such as

\[
\boxed{
\text{open},
\quad
\text{closed},
\quad
\text{compact},
\quad
\text{connected}.
}
\]

Measure theory studies concepts such as

\[
\boxed{
\text{measurable},
\quad
\text{null},
\quad
\text{finite measure},
\quad
\text{almost everywhere}.
}
\]

Neither notion of size replaces the other.

For example, the rational numbers are dense in \(\R\) but have Lebesgue
measure zero.

%%%%%%%%%%%%%%%%%%%%%%%%%%%%%%%%%%%%%%%%%%%%%%%%%%%%%%%%%%%%
\subsection{Measure and Probability}
\label{subsec:measure-and-probability}
%%%%%%%%%%%%%%%%%%%%%%%%%%%%%%%%%%%%%%%%%%%%%%%%%%%%%%%%%%%%

Modern probability theory is measure theory with total measure one.

The correspondence is

\[
\boxed{
\begin{array}{c|c}
\text{Measure Theory} & \text{Probability}\\
\hline
X & \Omega\\
\mathcal{M} & \mathcal{F}\\
\mu & \mathbb{P}\\
\text{measurable set} & \text{event}\\
\text{measurable function} & \text{random variable}
\end{array}
}
\]

This framework permits probability to handle continuous distributions,
infinite sample spaces, stochastic processes, and convergence of random
variables.

%%%%%%%%%%%%%%%%%%%%%%%%%%%%%%%%%%%%%%%%%%%%%%%%%%%%%%%%%%%%
\subsection{Measure and Integration}
\label{subsec:measure-and-integration}
%%%%%%%%%%%%%%%%%%%%%%%%%%%%%%%%%%%%%%%%%%%%%%%%%%%%%%%%%%%%

The Riemann integral partitions the domain into intervals.

Lebesgue integration instead organizes the theory around measurable sets
and their measures.

The basic progression is

\[
\boxed{
\text{indicator functions}
\longrightarrow
\text{simple functions}
\longrightarrow
\text{nonnegative measurable functions}
\longrightarrow
\text{general measurable functions}.
}
\]

This construction is developed in the next section.

%%%%%%%%%%%%%%%%%%%%%%%%%%%%%%%%%%%%%%%%%%%%%%%%%%%%%%%%%%%%
\subsection{Common Errors}
\label{subsec:measure-theory-common-errors}
%%%%%%%%%%%%%%%%%%%%%%%%%%%%%%%%%%%%%%%%%%%%%%%%%%%%%%%%%%%%

\subsubsection{Assuming Every Collection of Sets Is a Sigma-Algebra}

A sigma-algebra must be closed under complements and countable unions.

%%%%%%%%%%%%%%%%%%%%%%%%%%%%%%%%%%%%%%%%%%%%%%%%%%%%%%%%%%%%
\subsubsection{Confusing an Algebra with a Sigma-Algebra}

An algebra requires closure under finite unions.

A sigma-algebra requires closure under countable unions.

%%%%%%%%%%%%%%%%%%%%%%%%%%%%%%%%%%%%%%%%%%%%%%%%%%%%%%%%%%%%
\subsubsection{Confusing a Measurable Space with a Measure Space}

A measurable space is

\[
(X,\mathcal{M}).
\]

A measure space is

\[
(X,\mathcal{M},\mu).
\]

%%%%%%%%%%%%%%%%%%%%%%%%%%%%%%%%%%%%%%%%%%%%%%%%%%%%%%%%%%%%
\subsubsection{Assuming Measure Means Probability}

A general measure may satisfy

\[
\mu(X)>1
\]

or even

\[
\mu(X)=\infty.
\]

A probability measure specifically satisfies

\[
\mu(X)=1.
\]

%%%%%%%%%%%%%%%%%%%%%%%%%%%%%%%%%%%%%%%%%%%%%%%%%%%%%%%%%%%%
\subsubsection{Using Additivity on Overlapping Sets}

The formula

\[
\mu(A\cup B)
=
\mu(A)+\mu(B)
\]

requires disjointness.

In general,

\[
\boxed{
\mu(A\cup B)
=
\mu(A)+\mu(B)-\mu(A\cap B)
}
\]

when all quantities are finite.

%%%%%%%%%%%%%%%%%%%%%%%%%%%%%%%%%%%%%%%%%%%%%%%%%%%%%%%%%%%%
\subsubsection{Assuming Measure Zero Means Empty}

A null set may contain infinitely many or even uncountably many points.

%%%%%%%%%%%%%%%%%%%%%%%%%%%%%%%%%%%%%%%%%%%%%%%%%%%%%%%%%%%%
\subsubsection{Confusing Dense with Large Measure}

A dense set can have measure zero.

For example,

\[
\Q
\]

is dense in \(\R\) and has Lebesgue measure zero.

%%%%%%%%%%%%%%%%%%%%%%%%%%%%%%%%%%%%%%%%%%%%%%%%%%%%%%%%%%%%
\subsubsection{Assuming Every Subset of a Measurable Set Is Measurable}

This is false in general.

It is guaranteed for subsets of measurable null sets only when the measure
space is complete.

%%%%%%%%%%%%%%%%%%%%%%%%%%%%%%%%%%%%%%%%%%%%%%%%%%%%%%%%%%%%
\subsubsection{Confusing Borel and Lebesgue Measurability}

Every Borel set is Lebesgue measurable, but not every Lebesgue measurable
set is Borel.

%%%%%%%%%%%%%%%%%%%%%%%%%%%%%%%%%%%%%%%%%%%%%%%%%%%%%%%%%%%%
\subsubsection{Assuming Every Subset of \(\R\) Is Lebesgue Measurable}

Nonmeasurable sets exist.

%%%%%%%%%%%%%%%%%%%%%%%%%%%%%%%%%%%%%%%%%%%%%%%%%%%%%%%%%%%%
\subsubsection{Confusing Almost Everywhere with Everywhere}

A statement holding almost everywhere may fail on a nonempty null set.

%%%%%%%%%%%%%%%%%%%%%%%%%%%%%%%%%%%%%%%%%%%%%%%%%%%%%%%%%%%%
\subsubsection{Confusing Almost-Everywhere and In-Measure Convergence}

They are distinct notions.

On finite measure spaces,

\[
\text{a.e. convergence}
\Longrightarrow
\text{convergence in measure},
\]

but the reverse implication does not hold in general.

%%%%%%%%%%%%%%%%%%%%%%%%%%%%%%%%%%%%%%%%%%%%%%%%%%%%%%%%%%%%
\subsection{Applications}
\label{subsec:measure-theory-applications}
%%%%%%%%%%%%%%%%%%%%%%%%%%%%%%%%%%%%%%%%%%%%%%%%%%%%%%%%%%%%

\begin{applicationbox}
\textbf{Lebesgue Integration.}
Measure theory provides the sets and size structure required to define the
Lebesgue integral.

\medskip

\textbf{Probability Theory.}
Probability spaces are measure spaces whose total measure is one.

\medskip

\textbf{Functional Analysis.}
Spaces such as \(L^1\), \(L^2\), and \(L^p\) are defined using measurable
functions and Lebesgue integration.

\medskip

\textbf{Fourier Analysis.}
Modern Fourier theory relies heavily on \(L^p\) spaces and
measure-theoretic convergence.

\medskip

\textbf{Partial Differential Equations.}
Weak solutions and Sobolev spaces require functions defined and compared
up to sets of measure zero.

\medskip

\textbf{Dynamical Systems.}
Invariant measures describe long-term statistical behavior of dynamical
systems.

\medskip

\textbf{Ergodic Theory.}
Measure-preserving transformations connect dynamics, probability, and
statistical behavior.

\medskip

\textbf{Geometric Measure Theory.}
Measure is extended to curves, surfaces, fractals, and irregular geometric
objects.

\medskip

\textbf{Statistics.}
Probability distributions, densities, expectations, and likelihoods rely
on measure-theoretic foundations.
\end{applicationbox}

%%%%%%%%%%%%%%%%%%%%%%%%%%%%%%%%%%%%%%%%%%%%%%%%%%%%%%%%%%%%
\subsection{Exercises}
\label{subsec:measure-theory-exercises}
%%%%%%%%%%%%%%%%%%%%%%%%%%%%%%%%%%%%%%%%%%%%%%%%%%%%%%%%%%%%

\begin{exercise}[1]
Show that

\[
\{\varnothing,X\}
\]

is a sigma-algebra on \(X\).
\end{exercise}

\begin{exercise}[1]
Show that

\[
\mathcal{P}(X)
\]

is a sigma-algebra on \(X\).
\end{exercise}

\begin{exercise}[1]
If \(\mathcal{M}\) is a sigma-algebra, prove that

\[
\varnothing\in\mathcal{M}.
\]
\end{exercise}

\begin{exercise}[1]
Prove that a sigma-algebra is closed under countable intersections.
\end{exercise}

\begin{exercise}[1]
Explain the difference between

\[
(X,\mathcal{M})
\]

and

\[
(X,\mathcal{M},\mu).
\]
\end{exercise}

\begin{exercise}[1]
Determine the counting measure of

\[
\{2,4,6,8\}.
\]
\end{exercise}

\begin{exercise}[1]
For the Dirac measure \(\delta_0\), compute

\[
\delta_0([-1,1])
\]

and

\[
\delta_0((1,2)).
\]
\end{exercise}

\begin{exercise}[2]
Let

\[
X=\{1,2,3,4\}.
\]

Determine whether

\[
\mathcal{M}
=
\{
\varnothing,
\{1,2\},
\{3,4\},
X
\}
\]

is a sigma-algebra.
\end{exercise}

\begin{exercise}[2]
Prove monotonicity of measures:

\[
A\subseteq B
\Longrightarrow
\mu(A)\leq\mu(B).
\]
\end{exercise}

\begin{exercise}[2]
Prove countable subadditivity from countable additivity.
\end{exercise}

\begin{exercise}[2]
Prove continuity from below.
\end{exercise}

\begin{exercise}[2]
Prove continuity from above when

\[
\mu(A_1)<\infty.
\]
\end{exercise}

\begin{exercise}[2]
Compute the Lebesgue measure of

\[
(2,7],
\qquad
[0,3),
\qquad
\{1,2,3\}.
\]
\end{exercise}

\begin{exercise}[2]
Prove that every finite subset of \(\R\) has Lebesgue measure zero.
\end{exercise}

\begin{exercise}[2]
Prove that every countable subset of \(\R\) has Lebesgue measure zero.
\end{exercise}

\begin{exercise}[2]
Explain why

\[
m(\Q\cap[0,1])=0
\]

even though

\[
\Q\cap[0,1]
\]

is dense in \([0,1]\).
\end{exercise}

\begin{exercise}[2]
Show that if

\[
\mu(A)=0
\]

and \(B\subseteq A\) is measurable, then

\[
\mu(B)=0.
\]
\end{exercise}

\begin{exercise}[2]
Show that the indicator function

\[
\mathbf{1}_A
\]

is measurable if \(A\) is measurable.
\end{exercise}

\begin{exercise}[3]
Show that the intersection of any collection of sigma-algebras on \(X\)
is itself a sigma-algebra.
\end{exercise}

\begin{exercise}[3]
Use this fact to justify the existence of

\[
\sigma(\mathcal{C}).
\]
\end{exercise}

\begin{exercise}[3]
Prove that every closed subset of \(\R\) is Borel.
\end{exercise}

\begin{exercise}[3]
Prove that every countable subset of \(\R\) is Borel.
\end{exercise}

\begin{exercise}[3]
Prove that every continuous function

\[
f:\R\to\R
\]

is Borel measurable.
\end{exercise}

\begin{exercise}[3]
Suppose \(f\) and \(g\) are measurable. Prove that

\[
\max\{f,g\}
\]

and

\[
\min\{f,g\}
\]

are measurable.
\end{exercise}

\begin{exercise}[3]
Show that a pointwise limit of real-valued measurable functions is
measurable.
\end{exercise}

\begin{exercise}[3]
On \([0,1]\), analyze the almost-everywhere and in-measure convergence of

\[
f_n(x)=x^n.
\]
\end{exercise}

\begin{exercise}[3]
Construct an example showing that convergence in measure need not imply
pointwise convergence.
\end{exercise}

\begin{exercise}[3]
Prove the first Borel--Cantelli lemma.
\end{exercise}

\begin{exercise}[4]
Verify that Lebesgue outer measure is monotone.
\end{exercise}

\begin{exercise}[4]
Prove the countable subadditivity of Lebesgue outer measure.
\end{exercise}

\begin{exercise}[4]
Study the Carath\'eodory measurability criterion and prove that the
Carath\'eodory measurable sets form a sigma-algebra.
\end{exercise}

\begin{exercise}[4]
Explain why the construction of a Vitali set prevents a translation-
invariant countably additive extension of length to all subsets of
\(\R\).
\end{exercise}

\begin{exercise}[4]
Prove that on a finite measure space,

\[
f_n\to f
\text{ a.e.}
\Longrightarrow
f_n\to f
\text{ in measure}.
\]
\end{exercise}

\begin{exercise}[4]
Investigate the relationship between Lebesgue measurable sets and the
completion of the Borel sigma-algebra under Lebesgue measure.
\end{exercise}

%%%%%%%%%%%%%%%%%%%%%%%%%%%%%%%%%%%%%%%%%%%%%%%%%%%%%%%%%%%%
\subsection{Conceptual Summary}
\label{subsec:measure-theory-summary}
%%%%%%%%%%%%%%%%%%%%%%%%%%%%%%%%%%%%%%%%%%%%%%%%%%%%%%%%%%%%

Measure theory generalizes the concepts of length, area, volume, and
probability.

A sigma-algebra

\[
\boxed{
\mathcal{M}
}
\]

specifies which subsets are measurable.

A measurable space is

\[
\boxed{
(X,\mathcal{M}).
}
\]

A measure is a countably additive function

\[
\boxed{
\mu:\mathcal{M}\to[0,\infty].
}
\]

A measure space is

\[
\boxed{
(X,\mathcal{M},\mu).
}
\]

Countable additivity states that for pairwise disjoint measurable sets,

\[
\boxed{
\mu\left(
\bigcup_{n=1}^{\infty}A_n
\right)
=
\sum_{n=1}^{\infty}\mu(A_n).
}
\]

Measures satisfy monotonicity and countable subadditivity.

For increasing measurable sets,

\[
\boxed{
A_n\uparrow A
\Longrightarrow
\mu(A_n)\to\mu(A).
}
\]

For decreasing measurable sets with

\[
\mu(A_1)<\infty,
\]

\[
\boxed{
A_n\downarrow A
\Longrightarrow
\mu(A_n)\to\mu(A).
}
\]

Lebesgue measure extends ordinary geometric length:

\[
\boxed{
m([a,b])=b-a.
}
\]

Null sets satisfy

\[
\boxed{
\mu(N)=0.
}
\]

Statements holding outside a null set hold

\[
\boxed{
\text{almost everywhere}.
}
\]

Measurable functions preserve measurable structure through inverse images:

\[
\boxed{
B\text{ measurable}
\Longrightarrow
f^{-1}(B)\text{ measurable}.
}
\]

Outer measure and the Carath\'eodory criterion provide a systematic method
for constructing measures.

Measure theory therefore supplies the structural foundation for Lebesgue
integration, probability theory, \(L^p\) spaces, Fourier analysis,
functional analysis, partial differential equations, and modern
mathematical statistics.

%%%%%%%%%%%%%%%%%%%%%%%%%%%%%%%%%%%%%%%%%%%%%%%%%%%%%%%%%%%%
\subsection{Section Connections}
\label{subsec:measure-theory-connections}
%%%%%%%%%%%%%%%%%%%%%%%%%%%%%%%%%%%%%%%%%%%%%%%%%%%%%%%%%%%%

\chapterconnection{
Measure Theory extends Real Analysis by replacing elementary notions of
length, area, and volume with a countably additive theory of size.
Sigma-algebras provide the collection of sets on which measurement is
well behaved, while measures quantify those sets. Borel sigma-algebras
connect the theory directly to Point-Set Topology, and Lebesgue measure
extends ordinary geometric length while allowing substantially more
complicated sets. Null sets and almost-everywhere statements introduce a
new notion of analytical equivalence. Measurable functions and simple
functions provide the foundation for the Lebesgue integral developed in
the next section. The same framework becomes the mathematical foundation
of Probability Theory and later supports \(L^p\) spaces, Functional
Analysis, Fourier Analysis, Harmonic Analysis, Partial Differential
Equations, Ergodic Theory, and Geometric Measure Theory.
}

%%%%%%%%%%%%%%%%%%%%%%%%%%%%%%%%%%%%%%%%%%%%%%%%%%%%%%%%%%%%
% SECTION SOURCE TRACKING
%%%%%%%%%%%%%%%%%%%%%%%%%%%%%%%%%%%%%%%%%%%%%%%%%%%%%%%%%%%%

%------------------------------------------------------------
% 5.6 Measure Theory
%------------------------------------------------------------
%
% Source basis:
%   - Standard undergraduate and beginning graduate measure theory.
%   - Standard construction and elementary properties of
%     Lebesgue measure.
%
% Used for:
%   - algebras of sets;
%   - sigma-algebras;
%   - measurable spaces;
%   - generated sigma-algebras;
%   - Borel sigma-algebras;
%   - measures and measure spaces;
%   - finite and sigma-finite measures;
%   - probability measures;
%   - counting measure;
%   - Dirac measure;
%   - monotonicity;
%   - countable subadditivity;
%   - continuity from below;
%   - continuity from above;
%   - outer measure;
%   - Caratheodory measurability;
%   - Lebesgue measure;
%   - translation invariance;
%   - null sets;
%   - complete measures;
%   - Borel versus Lebesgue measurability;
%   - nonmeasurable sets;
%   - measurable functions;
%   - indicator functions;
%   - simple functions;
%   - pointwise limits of measurable functions;
%   - almost-everywhere convergence;
%   - convergence in measure;
%   - limsup and liminf of sets;
%   - Borel--Cantelli preview;
%   - product measure preview;
%   - higher-dimensional Lebesgue measure;
%   - regularity of Lebesgue measure.
%
% Notes:
%   - The rigorous theory of sequences, limits, continuity, and
%     Riemann integration is developed in Section 5.5.
%   - The construction and convergence theory of the Lebesgue
%     integral is developed in Section 5.7.
%   - General topological concepts and Borel constructions connect
%     to Chapter 4.
%   - Probability measures are developed extensively in Chapter 7.
%   - Product measures and Fubini--Tonelli theory are expanded in
%     the Lebesgue Integration section.
%
%%%%%%%%%%%%%%%%%%%%%%%%%%%%%%%%%%%%%%%%%%%%%%%%%%%%%%%%%%%%